\section{Threats to Validity}
\label{sec:threats_to_validity}

This section addresses the potential threats to validity across different dimensions and the strategies adopted to mitigate them \cite{book:wohlin}.

\textbf{Internal Validity.}
For internal validity, factors such as hardware and software influence on the stability of the results were considered. To mitigate these interferences, all tests were executed in a local environment, on the same machine and operating system. Additionally, each experiment was run ten times, ensuring that the observed results are consistent and not merely due to random chance or transient system variations. This repetition enhances the reliability of the findings by providing more robust and representative data averages.

\novo{Uma ameaça adicional à validade interna nesta versão estendida diz respeito à assimetria do design experimental: a minificação foi aplicada exclusivamente às respostas CSR (formato JSON), enquanto as respostas SSR (HTML pré-renderizado) não foram submetidas a minificação. Essa decisão foi deliberada — a minificação de HTML envolve um mecanismo técnico distinto da minificação de JSON e está fora do escopo deste estudo — mas limita a possibilidade de comparações simétricas entre os efeitos da minificação nas duas estratégias de renderização. As conclusões sobre os efeitos da minificação devem, portanto, ser interpretadas como específicas para arquiteturas CSR. Adicionalmente, a compressão gzip introduz uma sobrecarga de CPU no lado do servidor que não foi mensurada neste estudo; as melhorias de latência observadas podem não se generalizar integralmente para ambientes onde a capacidade de CPU do servidor é limitada.}

\textbf{External Validity.}
Regarding external validity, a recognized limitation of this study is that the experiments were conducted within a single execution environment. All experiments were executed on a single hardware configuration (a MacBook Air M2) and within a single browser environment (Chromium). This restricted setup constrains the generalizability of the findings, as performance and behavior may vary across distinct hardware platforms, operating systems, and browsers employing different rendering engines. However, with the environment used, relevant insights were found for the study proposal, which can be expanded in future work.

\novo{Uma ameaça adicional à validade externa diz respeito às condições de rede. Todos os experimentos foram conduzidos em ambiente localhost, sem latência de rede simulada nem limitação de largura de banda. Em cenários de implantação reais, a compressão gzip teria um efeito ainda mais pronunciado, pois reduz o volume de bytes transmitidos em conexões onde a largura de banda é fator limitante. Os valores absolutos de LCP e TTI reportados neste estudo podem, portanto, subestimar o benefício do gzip em ambientes práticos com restrições de rede.~\citet{ollila:2021} observou que as diferenças de desempenho entre estratégias de renderização variam em função do motor JavaScript utilizado, sugerindo que a magnitude das diferenças observadas pode divergir sob implementações alternativas de frameworks como Next.js, Nuxt ou SvelteKit. Por fim, o domínio de aplicação — uma galeria de fotos servindo metadados de imagens — representa uma classe específica de aplicações web de conteúdo estático; a generalização dos achados para SPAs altamente interativas, dashboards ou plataformas de e-commerce é limitada, conforme apontado por~\citet{gangishetti:2026} em seu estudo sobre plataformas de conteúdo em larga escala.}

\textbf{Conclusion Validity.}
With respect to conclusion validity, several techniques exist to improve application performance, such as data minification and compression, asynchronous loading (lazy loading), the use of Content Delivery Networks (CDNs), image preloading, and data caching. \novo{Esta versão estendida investiga diretamente duas dessas técnicas — minificação e gzip — em um experimento fatorial controlado, endereçando uma limitação explicitamente identificada no estudo original.} However, the objective of this study was to analyze the impact of the rendering approach (SSR and CSR) and data size on the user experience. To ensure that the evaluation focus remained on rendering performance, the application used only HTML, CSS, and JavaScript without any external libraries.

\novo{Em relação à validade estatística de conclusão, a análise estendida emprega o teste post-hoc de Dunn com correção de Holm para controlar a taxa de erro familiar nas múltiplas comparações par a par, reduzindo o risco de erros do Tipo I~\cite{book:wohlin}. Adicionalmente, o uso do Delta de Cliff com intervalos de confiança via bootstrap fornece estimativas de tamanho de efeito robustas a pressupostos distribucionais e a amostras pequenas~\cite{efron:1994,Kitchenham2024}. Contudo, com $n = 10$ observações por célula, o poder estatístico para detectar efeitos pequenos ou médios ($|\delta| < 0{,}33$) pode ser insuficiente; todas as conclusões reportadas limitam-se, portanto, aos efeitos grandes observados.}

\novo{\textbf{Construct Validity.}
A utilização das métricas Core Web Vitals do Google (LCP, FCP, SI, TTI, CLS, FID e Lighthouse Score) como indicadores de qualidade de experiência do usuário é fundamentada no framework estabelecido pelo Google e corroborada por pesquisas que demonstram correlação entre essas métricas e indicadores de comportamento do usuário, como taxa de rejeição, duração de sessão e taxas de conversão~\cite{sikder:2016}. Contudo, o Lighthouse Score é uma métrica composta sintética que pode não capturar todas as dimensões da experiência do usuário percebida, especialmente em aplicações com padrões de interação complexos. Adicionalmente, o mapeamento de \texttt{format=html} para SSR e \texttt{format=json} para CSR representa uma implementação simplificada de cada paradigma; frameworks SSR modernos (como Next.js e Nuxt) e frameworks CSR (como React, Vue e Svelte) podem apresentar características de desempenho distintas, conforme documentado por~\citet{ollila:2021}. Os achados devem ser compreendidos como caracterização das compensações (\textit{trade-offs}) fundamentais entre a transferência de HTML pré-renderizado e a renderização de JSON no lado do cliente, e não como avaliação de frameworks específicos de produção.}
